\section{Introduction}

In this paper, we study the modules over face rings, introduced by Athanasiadis and Stanley, whose Hilbert functions are the relative local $h$-vectors of quasi-geometric homology triangulations of simplices, a broad class of formal subdivisions that includes all geometric triangulations and is natural from the point of view of combinatorial commutative algebra.  See Section 2.1 for the precise definition and further references.

Fix an infinite field $k$.  Let $\sigma \colon \Gamma \to 2^V$ be a quasi-geometric homology triangulation of a simplex, and let $E$ be a face of $\Gamma$. Say that a face $G \in \Gamma$ is \emph{interior} if $\sigma(G) = V$, and let $I$ be the ideal in the face ring $k[\lk_\Gamma(E)]$ generated by the faces that are interior relative to $E$, i.e. 
\[
 I = (x^F : F \sqcup E \mbox{ is interior} ).
\]
Let $d = |V|-|E|$, which is the Krull dimension of $k[\lk_\Gamma(E)]$, and let $\theta_1, \ldots, \theta_{d}$ be a special l.s.o.p., as in \cite{Stanley92, Athanasiadis12b}. See also Section~\ref{ss:sr}, where we recall the definition and construction of special l.s.o.p.s.

\begin{definition}
The \emph{local face module} $L(\Gamma,E)$ is defined as the image of $I$ in $k[\lk_\Gamma(E)]/(\theta_1, \ldots, \theta_d)$.
\end{definition}

Note that $L(\Gamma,E)$ is a finite dimensional graded $k$-vector space. The \emph{local $h$-vector} is its Hilbert function: 
\[
\ell(\Gamma, E)  \coloneqq  (\ell_0, \ldots, \ell_{d}), \quad \quad \mbox{ where } \ \ell_i  \coloneqq  \dim L(\Gamma,E)_i.
\]
The local face module $L(\Gamma,E)$ depends on the choice of a special l.s.o.p., but $\ell(\Gamma,E)$ is an invariant of the triangulation with the symmetry $\ell_i = \ell_{d- i}$. See Section~\ref{sec:triangulations} for details and references. In the past few years, there has been significant research activity on the combinatorics of local $h$-vectors and relations to intersection homology \cite{Athanasiadis16, KatzStapledon16, Stapledon17, deCataldoMiglioriniMustata18}.  Recent advances include a proof that every non-negative integer vector satisfying $\ell_0 = 0$ and $\ell_i = \ell_{d-i}$ is the local $h$-vector of a quasi-geometric triangulation for $E = \emptyset$ \cite{JKMS}, and a relative hard Lefschetz theorem that yields unimodality of local $h$-vectors for regular subdivisions in a more general setting (for regular nonsimplicial polyhedral subdivisions that are not necessarily rational) \cite{Karu19}.

 

Here, we investigate the local face modules $L(\Gamma, E)$ using methods of combinatorial commutative algebra. In particular, we describe natural combinatorial resolutions of these modules as well as natural maps of $k[\lk_\Gamma(E)]$-modules, $L(\Gamma,E) \to L(\Gamma, E')$, for $E \subset E'$.  Our first theorem gives explicit generators for the kernel of the natural map $I \to k[\lk_\Gamma(E)]/(\theta_1, \ldots, \theta_d)$. Moreover, we extend this to an exact sequence of graded $k[\lk_\Gamma(E)]$-modules in which each term is a direct sum of degree-shifted monomial ideals.

Label the vertices of the simplex $V = \{ v_1, \ldots, v_n \}$.  For a subset $U \subset V$, let $U^c  \coloneqq  V \smallsetminus U$. After relabeling, we may assume that $\sigma(E)^c = \{v_1, \ldots, v_b\}$. Given $S \subset \{v_1, \ldots, v_d\}$, 
we define the ideal $I_S \subset k[\lk_\Gamma(E)]$ by
\begin{equation*}
I_S  \coloneqq  ( x^F : \, \sigma(F \sqcup E)^c \subset S). 
\end{equation*}
Note that $I_{S'} \subset I_{S}$ for $S' \subset S$, and $I_S$ depends only on $S \cap \{v_1, \ldots, v_b\}$. For instance, $I_{\emptyset} = I$ and $I_{S} = k[\lk_{\Gamma}(E)]$ if $\{v_1, \ldots, v_b \} \subset S$. By the definition of a special l.s.o.p. (Definition~\ref{def:special}), after reordering, we may assume
\[
\supp(\theta_i) \subset \{ w \in \lk_\Gamma(E) : v_i \in \sigma(w) \}, 
\]
for $1 \leq i \leq b$.  As a consequence, for any $v_i \in S$, multiplication by $\theta_i$ induces a degree 1 map $\lambda_i \colon I_S \to I_{S \smallsetminus \{v_i\}}$.

\begin{theorem}\label{thm:resolution}
There is an exact sequence of graded $k[\lk_{\Gamma}(E)]$-modules
 $$0  \to k[\lk_{\Gamma}(E)][-d] \to \bigoplus_{ |S|  = d - 1 } I_{S}[-(d-1)] \to \dotsb  \to \bigoplus_{|S| = 1} I_{S}[-1] \to I \to L(\Gamma,E) \to 0,$$
where,  
for $S= \{v_{i_0}, \ldots, v_{i_k}\}$, with $i_0 < \cdots < i_k$, the differential restricted to $I_S$ is $\oplus_{j = 0}^k (-1)^j \lambda_{i_j}$.
\end{theorem}

\begin{corollary}\label{cor:presentation}
The kernel of the surjection $I \to L(\Gamma, E)$ is the ideal $J$ generated by 
 $$ \left \{ \theta_i \cdot x^{F} : F \sqcup E \mbox{ is interior } \right \}  \cup  \left \{ \theta_{j} \cdot x^G : \sigma(G \sqcup E) = \{v_j\}^c, \mbox{ for } 1 \leq j \leq b \right \}.$$
\end{corollary}
